\section{Προκλήσεις και Ανοικτά Ερευνητικά Προβλήματα}

Παρά τη σημαντική πρόοδο στον συνδυασμό \en{Large Language Models (LLMs)}, \en{LoRA} και \en{Federated Learning}, εξακολουθούν να υπάρχουν αρκετές προκλήσεις και ανοικτά ερευνητικά ζητήματα που επηρεάζουν την πρακτική εφαρμογή αυτών των τεχνολογιών, ιδιαίτερα σε πολύπλοκα και μεγάλης κλίμακας περιβάλλοντα όπως τα δίκτυα \en{5G} και \en{6G}. Ένας από τους σημαντικότερους περιορισμούς είναι το επικοινωνιακό κόστος, καθώς παρότι τεχνικές όπως το \en{LoRA} μειώνουν το μέγεθος των ενημερώσεων του μοντέλου, η συχνή ανταλλαγή πληροφοριών και ο μεγάλος αριθμός συμμετεχόντων πελατών μπορούν να προκαλέσουν συμφόρηση δικτύου και αυξημένη καθυστέρηση \cite{kairouz2021advances}. Επιπλέον, τα συστήματα \en{Federated Learning} αντιμετωπίζουν συχνά προβλήματα σύγκλισης λόγω ασύγχρονων ενημερώσεων, μεταβαλλόμενης συμμετοχής των πελατών και διαφοροποιημένων συνθηκών εκπαίδευσης, γεγονός που καθιστά δύσκολη τη διασφάλιση σταθερής και αποδοτικής εκπαίδευσης. Παράλληλα, η ετερογένεια δεδομένων αποτελεί σημαντική πρόκληση, καθώς τα δεδομένα σε ομοσπονδιακά περιβάλλοντα είναι συνήθως μη ανεξάρτητα και μη ομοιόμορφα κατανεμημένα (\en{non-IID}), επηρεάζοντας αρνητικά την απόδοση των μοντέλων λόγω των διαφοροποιημένων τοπικών κατανομών μεταξύ των πελατών. Επιπρόσθετα, η επεκτασιμότητα των συστημάτων σε περιβάλλοντα \en{5G}/\en{6G} απαιτεί αποδοτική διαχείριση πόρων, προσαρμοστικές στρατηγικές επικοινωνίας και ανθεκτικό σχεδιασμό, εξαιτίας του μεγάλου αριθμού συσκευών και των συνεχώς μεταβαλλόμενων συνθηκών λειτουργίας. Τέλος, παρά την ενίσχυση της ιδιωτικότητας που προσφέρει το \en{Federated Learning}, ζητήματα όπως η διαρροή \en{gradients}, οι επιθέσεις αντιστροφής μοντέλου (\en{model inversion}) και οι προκαταλήψεις στα δεδομένα εξακολουθούν να αποτελούν σημαντικές απειλές \cite{zhu2019deep}. Παράλληλα, ζητήματα ηθικής όπως η δικαιοσύνη, η διαφάνεια και η λογοδοσία παραμένουν κρίσιμοι παράγοντες για την ασφαλή και υπεύθυνη ανάπτυξη αυτών των συστημάτων.